%! TEX root = ../main.tex

\newpage

\section{Giới thiệu}

\subsection{Mục đích}

<Xác định mục tiêu của tài liệu này và đối tượng dự định đọc nó. (VD: Tài liệu thiết kế phần mềm 
này mô tả thiết kế kiến trúc và thiết kế chi tiết của XX…).> 

\subsection{Phạm vi}

<Viết mô tả và phạm vi của phần mềm và giải thích các lợi ích, mục đích và mục tiêu của dự án.> 

\subsection{Bảng chú giải thuật ngữ}
\begin{itemize}
  \item Đây là mục tùy chọn 
  \item Đây là mục tùy chọn 
Định nghĩa các từ viết tắt, các thuật ngữ được sử dụng trong tài liệu mà chúng gần như 
không được biết đến bởi người đọc.

\noindent
\begin{tabularx}{\textwidth}{|>{\raggedleft\arraybackslash}p{1cm}|p{5cm}|X|}
    \hline
    \rowcolor{gray!25} 
    % \multicolumn{1}{|c|}{...} ép riêng ô STT nằm căn giữa
    \multicolumn{1}{|c|}{\textbf{STT}} & \textbf{Thuật ngữ/Từ viết tắt} & \textbf{Định nghĩa / Giải thích} \\ \hline
    1 & PM & Project Manager \\ \hline
    2 & QA & Quality Assurance \\ \hline
    3 & UX & User Experience \\ \hline
\end{tabularx}
  
\end{itemize}

\subsection{Tài liệu tham khảo}

\begin{itemize}
    \item Đây là mục tùy chọn
    \item Đây là mục tùy chọn 
Liệt kê ra bất cứ tài liệu hay địa chỉ website nào mà tài liệu này tham khảo tới. Cung cấp đủ 
thông tin để người đọc có thể tìm bản sao của từng tài liệu tham khảo, bao gồm: tiêu đề, 
tác giả, số phát hành, ngày, nguồn hay nơi cung cấp. > 
\end{itemize}

\subsection{Tổng quan về tài liệu}
<Cung cấp một cái nhìn tổng quan về tài liệu này và sự tổ chức của nó.>
